% Section 4, printed pp. 26--35.
\section{The Picard scheme}\label{kps-sec-scheme}

\begin{definition}[Picard scheme]\label{kps-def-picard-scheme}
\dcref{4.1}\source{kleiman-picard:page-0026}
If one of the relative Picard functors is representable, its representing
scheme is the Picard scheme $\IPic_{X/S}$. Representability of the etale
sheaf implies representability of the fppf sheaf; any representing functor is
therefore uniquely the same scheme.
\end{definition}

\begin{definition}[Abel map]\label{kps-def-abel-map}
\dcref{4.6}\source{kleiman-picard:page-0027}
When $\IPic_{X/S}$ exists, the universal line bundle (when available) sends a
relative effective divisor to its class, defining the Abel map
$A_{X/S}:\Div_{X/S}\to\IPic_{X/S}$.
\end{definition}

\begin{theorem}[Main Picard representability theorem]\label{kps-thm-main}
\dcref{4.8}\source{kleiman-picard:page-0027}
Assume $f:X\to S$ is projective Zariski locally over $S$, flat, and has
integral geometric fibers. Then $\IPic_{X/S}$ exists, is separated and locally
of finite type over $S$, and represents $\Pic_{(X/S)\etale}$. If $S$ is
Noetherian and $X/S$ is projective, it is a disjoint union of open subschemes,
each an increasing union of open quasi-projective $S$-schemes.
\uses{kps-def-picard-scheme,kps-def-relative,kps-def-abel-map,kps-thm-divisor-representability,kps-thm-linear-systems,kps-lem-quotient}
\end{theorem}
\begin{proof}
The assertion is local on $S$, so first reduce to the Noetherian projective
case. Write $P=\Pic_{(X/S)\etale}$. For a polynomial $\phi$, let $P^\phi$
be the etale subsheaf of classes represented by $\LL$ with
\[
  \chi(X_t,\LL_t^{-1}(n))=\phi(n)\quad\text{for every }t\text{ and }n.
\]
This condition is invariant under etale base change by flat base change for
fiber cohomology. Given $T\to P$ represented by $\LL'$ on an etale cover
$T'\to T$, the subset
\[
 T'^\phi=\{t'\mid \chi(X_{t'},\LL_{t'}'^{-1}(n))=\phi(n)\text{ for all }n\}
\]
is open by semicontinuity. Isomorphic pullbacks on an etale cover of
$T'\times_TT'$ show $T'^\phi$ is the inverse image of an open $T^\phi\subset T$.
The same descent argument shows that $T^\phi$ represents
$P^\phi\times_PT$. Thus the $P^\phi$ are disjoint open subfunctors covering
$P$.

For $m\in\mathbb Z$, let $P_m^\phi$ impose in addition
\[
 R^if_{T*}(\LL(n))=0\quad(i\ge1,\ n\ge m).
\]
By the cohomology-and-base-change criterion this is equivalent to
\[
 H^i(X_t,\LL_t(n))=0\quad(i\ge1,\ n\ge m,\ t\in T).
\]
Indeed, fiberwise vanishing implies the direct-image vanishing; conversely,
descending induction on $i$, starting with Serre vanishing, uses the long exact
sequence after tensoring by an arbitrary quasi-coherent module and then by
$k(t)$ to recover the fiber vanishing. Flat base change preserves the latter
condition, so $P_m^\phi$ is well defined. Serre's finiteness theorem makes the
vanishing locus open, and the $P_m^\phi$ form a nested cover of $P^\phi$.

Put $\phi_0(n)=\phi(m+n)$. Tensoring by $\mathcal O_X(m)$ gives an isomorphism
$P_m^\phi\simeq P_0^{\phi_0}$. It remains to represent the latter. After
replacing $S$ by an open-and-closed piece, assume
$\psi(n)=\chi(X_s,\mathcal O_{X_s}(n))$ is constant. Form
\[
 Z=P_0^{\phi_0}\times_P\Div_{X/S},\qquad
 \theta(n)=\psi(n)-\phi_0(n).
\]
By representability of relative divisors, $Z$ is an open subscheme of
$\Div_{X/S}$ and lies in the projective Hilbert-polynomial stratum
$\Hilb_{X/S}^{\theta}$; hence it is quasi-projective.

We check that $\alpha:Z\to P_0^{\phi_0}$ is an etale-sheaf surjection. Take
$\lambda\in P_0^{\phi_0}(T)$, represented by $\LL'$ on an etale cover
$T'\to T$. The pullback $T'\times_{P_0^{\phi_0}}Z$ is the linear-system
functor for $\LL'$, hence is $\PP(\QQ)$ by Theorem~3.13. Here
$H^1(X_t,\LL'_t)=0$, so the module $\QQ$ is locally free. Therefore
$\PP(\QQ)\to T'$ is smooth and has an etale-local section, producing an
etale cover $T_1\to T$ and a lift $T_1\to Z$ of $\lambda$. Thus $\alpha$ is
surjective. Taking $T=Z$ in the same argument shows that
$R=Z\times_{P_0^{\phi_0}}Z$ is a scheme and that $R\to Z$ is smooth and
proper. Lemma~4.9 now represents $P_0^{\phi_0}$ by a quasi-projective scheme.
Consequently every $P_m^\phi$ is an increasing union of open quasi-projective
schemes, and their disjoint union represents $P$. The construction also gives
separatedness and local finite type, proving both assertions.
\end{proof}

\begin{lemma}[Quotient lemma]\label{kps-lem-quotient}
\dcref{4.9}\source{kleiman-picard:page-0030}
Let $\alpha:Z\to P$ be a surjection of etale sheaves and
$R=Z\times_P Z$. If $Z$ is a quasi-projective $S$-scheme, $R$ is representable,
and $R\to Z$ is smooth and proper, then $P$ is represented by a
quasi-projective $S$-scheme and $\alpha$ is smooth.
\end{lemma}
\begin{proof}
The map $R\to Z\times_S Z$ is a proper monomorphism, hence a closed
immersion, and defines a flat proper equivalence relation. The Hilbert-scheme
construction of the quotient gives a faithfully flat projective map $Z\to Q$
with $R=Z\times_QZ$. Fiberwise smoothness of $R\to Z$ implies $Z\to Q$ is
smooth. It is an etale-sheaf surjection, so both $Q$ and $P$ are the
coequalizer of $R\rightrightarrows Z$ and are uniquely isomorphic.
\end{proof}

\begin{remark}[Quotient criterion]\label{kps-ex-quotient-criterion}
\source{kleiman-picard:page-0031}
A map of etale sheaves is a surjection exactly when its target is the
coequalizer of $F\times_G F\rightrightarrows F$.
\end{remark}
\begin{remark}[Finite type subschemes]\label{kps-ex-quasi-projective}
\source{kleiman-picard:page-0031}
Every finite-type subscheme of $\IPic_{X/S}$ under the hypotheses of Theorem
4.8 is quasi-projective.
\end{remark}
\begin{remark}[Abel map properness]\label{kps-ex-abel-proper}
\source{kleiman-picard:page-0031}
If a Poincare sheaf exists, the Abel map is projective Zariski locally; in
general it is proper after base change and descent.
\end{remark}
\begin{remark}[Curve Abel embedding]\label{kps-ex-curve-abel}
\source{kleiman-picard:page-0031}
For a flat locally projective family of integral curves of genus at least one,
the smooth locus embeds naturally as a closed subscheme in the Picard scheme.
\end{remark}

\begin{example}[Mumford's nonintegral-fiber obstruction]\label{kps-ex-mumford}
\dcref{4.14}\source{kleiman-picard:page-0031}
For $S=\Spec\mathbb R[[t]]$ and the conic $x^2+y^2=t$, the special geometric
fiber splits into conjugate lines. After extending to $\mathbb C[[t]]$, the
relative Picard functor is a disjoint union of nonseparated schemes obtained by
gluing copies of the base at the closed point, indexed by multidegrees. The
conjugate branches cannot lie in one affine open, so descent would contradict
affineness; consequently the Picard functor over $S$ is not representable.
\end{example}

\begin{proposition}[Local finite type]\label{kps-prop-lft}
\dcref{4.17}\source{kleiman-picard:page-0033}
If $\IPic_{X/S}$ exists and represents the fppf Picard sheaf, then it is locally
of finite type over $S$.
\end{proposition}
\begin{proof}
The Hilbert-polynomial decomposition in the proof of Theorem 4.8 is locally
finite type, and the open pieces cover the Picard scheme.
\end{proof}

\begin{theorem}[Mumford's reduced-fiber representability]\label{kps-thm-mumford}
\dcref{4.18.1}\source{kleiman-picard:page-0033}
If $X/S$ is projective and flat with reduced connected geometric fibers, and the
irreducible components of ordinary fibers are geometrically irreducible, then
the Picard scheme exists.
\notready
\end{theorem}

\begin{theorem}[Generic representability]\label{kps-thm-relative-representability}
\dcref{4.18.2}\source{kleiman-picard:page-0034}
If $X/S$ is proper and $S$ is integral, there is a nonempty open $V\subset S$
such that $\IPic_{X_V/V}$ exists, represents the fppf Picard functor, and is a
disjoint union of open quasi-projective subschemes.
\notready
\end{theorem}

\begin{corollary}[Complete varieties]\label{kps-cor-complete-representability}
\dcref{4.18.3}\source{kleiman-picard:page-0034}
If $X$ is complete over a field, then $\IPic_{X/k}$ exists, represents the fppf
Picard functor, and is a disjoint union of open quasi-projective subschemes.
\uses{kps-thm-relative-representability}
\end{corollary}

\begin{theorem}[Flattening stratification]\label{kps-thm-flattening}
\dcref{4.18.4}\source{kleiman-picard:page-0035}
For proper $X/S$ and a coherent sheaf $\mathcal F$, the functor of $S$-schemes
over which $\mathcal F$ becomes flat is represented by an unramified finite-type
$S$-scheme.
\notready
\end{theorem}

\begin{theorem}[Relative representability map]\label{kps-thm-relative-map}
\dcref{4.18.5}\source{kleiman-picard:page-0035}
If $S$ is integral and $X'\to X$ is a surjective map of proper $S$-schemes,
then over a nonempty open $V\subset S$ the induced map of fppf Picard functors
is representable by quasi-affine maps of finite type.
\notready
\end{theorem}

\begin{theorem}[Artin's algebraic-space representability]\label{kps-thm-algebraic-space}
\dcref{4.18.6}\source{kleiman-picard:page-0036}
For a flat proper finitely presented map of algebraic spaces with formation of
$f_*\mathcal O_X$ commuting with base change, the Picard functor exists as an
algebraic space locally of finite presentation over $S$.
\notready
\end{theorem}

\begin{remark}[Flattening and algebraic spaces]\label{kps-rmk-existence}
\source{kleiman-picard:page-0035}
The relative representability theorem follows from flattening stratifications
and Chow's lemma. The algebraic-space version removes projectivity; its proof
uses the same Hilbert and quotient constructions.
\end{remark}
